\lecture{13}{5. December 2025}{Neumann and Mixed Boundary Conditions. Irregular Boundary.}

\begin{problemwithimage}{p13_1}{13.1}
  Consider the two-dimensional Laplace equation
  \begin{equation*}
    u_{x x}(x,y) + u_{y y}(x,y) = 0
  \end{equation*}
  on the region $R$ with mixed boundary conditions as depicted on the figure.
  Find the corresponding system of approximate difference equations for mesh size $h = 3$.
\end{problemwithimage}

\begin{derivation}
  The corresponding difference equation is (L. 12 p. 5):
  \begin{equation*}
    u_{i + 1,j} + u_{i,j+1} + u_{i-1,j} + u_{i,j-1} - 4u_{i,j} = 0
  .\end{equation*}
  Our boundary meshpoints are
  \begin{align*}
    u_{00} &= 100, & u_{10} &= 100, & u_{20} &= 100, & u_{30} &= 100, \\
    u_{01} &= 100, & u_{31} &= 100, & u_{02} &= 100, & u_{32} &= 100, \\
           &&  u_{03} &= 100, & u_{33} &= 100. &&
  \end{align*}

  The inner meshpoint $P_{11}$ is found as
  \begin{align*}
    u_{21} + u_{12} + u_{01} + u_{10} - 4u_{11} &= 0 \\
    u_{21} + u_{12} + 100 + 100 - 4u_{11} &= 0 \\
    -4u_{11} + u_{21} + u_{12} &= -200
  .\end{align*}
  
  Now, the inner meshpoint $P_{21}$ is found as
  \begin{align*}
    u_{31} + u_{22} + u_{11} + u_{20} - 4u_{21} &= 0 \\
    100 + u_{22} + u_{11} + 100 - 4u_{21} &= 0 \\
    u_{11} - 4u_{21} + u_{22} &= -200
  .\end{align*}

  Now $P_{12}$ is found as
  \begin{align*}
    u_{22} + u_{13} + u_{02} + u_{11} - 4u_{12} &= 0 \\
    u_{22} + u_{31} + 100 + 100 - 4u_{12} &= 0 \\
    u_{11} -4u_{12} + u_{22} + u_{31} &= -200
  .\end{align*}
  
  Now we find $P_{22}$ as
  \begin{align*}
    u_{32} + u_{23} + u_{12} + u_{21} - 4u_{22} &= 0 \\
    100 + u_{23} + u_{12} + u_{21} - 4u_{22} &= 0 \\
    u_{21} + u_{12} - 4u_{22} + u_{23} &= -100
  .\end{align*}

  Now we move on to the boundary point $P_{13}$, which we find as
  \begin{align*}
    u_{23} + u_{14} + u_{03} + u_{12} - 4u_{13} &= 0 \\
    u_{23} + u_{14} + 100 + u_{12} - 4u_{13} &= 0 \\
    u_{12} - 4u_{13}+ u_{23} + u_{14} &= -100
  .\end{align*}
  To find an expression for $u_{14}$, we use the directional derivative as
  \begin{align*}
    118 = u_n(P_{13}) = \frac{\partial u_{13}}{\partial y} &\approx \frac{u_{14} - u_{12}}{2h} = \frac{u_{14} - u_{12}}{6} \\
    \implies u_{14} &= u_{12} + 708
  .\end{align*}
  We insert this into the previous expression to get
  \begin{align*}
    u_{12} - 4u_{13} + u_{23} + u_{12} + 708 &= -100 \\
    2u_{12} - 4u_{13} + u_{23} &= -808
  .\end{align*}

  Lastly we find the boundary point $P_{23}$ as
  \begin{align*}
    u_{33} + u_{24} + u_{13} + u_{22} - 4u_{23} &= 0 \\
    100 + u_{24} + u_{13} + u_{22} - 4u_{23} &= 0 \\
    u_{13} + u_{22} - 4u_{23} + u_{24} &= -100
  .\end{align*}
  We, once again, use the directional derivative as
  \begin{align*}
    118 = u_n(P_{23}) = \frac{\partial u_{23}}{\partial y} &\approx \frac{u_{24} - u_{22}}{2h} = \frac{u_{24} - u_{22}}{6} \\
    u_{24} &= u_{22} + 708
  .\end{align*}
  And by insertion we get
  \begin{align*}
    u_{13} + u_{22} - 4u_{23} + u_{22} + 708 &= -100 \\
    2u_{22} + u_{13} - 4u_{23} &= -808
  .\end{align*}
\end{derivation}

\finalresult{
\begin{aligned}
  -4u_{11} + u_{21} + u_{12} &= -200 \\
  u_{11} - 4u_{21} + u_{22} &= -200 \\
  u_{11} - 4u_{12} + u_{22} + u_{31} &= -100 \\
  u_{21} + u_{12} - 4u_{22} + u_{23} &= -100 \\
  2u_{12} - 4u_{13} + u_{23} &= -808 \\
  2u_{22} + u_{13} - 4u_{23} &= -808
.\end{aligned}
}

\begin{problemwithimage}{p13_1}{13.2}
  Consider the two-dimensional Poisson equation
  \begin{equation*}
    u_{x x}(x,y) + u_{yy}(x,y) = x + y
  \end{equation*}
  on the region $R$ with mixed boundary conditions as depicted on the figure.
  Find the corresponding system of approximate solutions for mesh size $h = \num{4,5}$.
\end{problemwithimage}

\begin{derivation}
  We have the difference equation in mesh notation:
  \begin{equation*}
    u_{i+1,j} + u_{i,j+1} + u_{i-1,j} + u_{i,j-1} - 4u_{ij} = h^2(ih + jh)
  .\end{equation*}
  The known boundary mesh points are
  \begin{align*}
    u_{02} &= 100, & u_{10} &= 100, & u_{20} &= 100, \\
    u_{01} &= 100, & u_{21} &= 100, & u_{22} &= 100
  .\end{align*}

  Now, we find the inner meshpoint $P_{11}$ as
  \begin{align*}
    u_{21} + u_{12} + u_{01} + u_{10} - 4u_{11} &= \num{4,5}^2 \cdot \left( 1 \cdot \num{4,5} + 1 \cdot \num{4,5}  \right) \\
      100 + u_{12} + 100 + 100 - 4u_{11} &= \num{182,25} \\
      -4u_{11} + u_{12} &= \num{-117,75} 
  .\end{align*}

  The boundary point $P_{12}$ is found as
  \begin{align*}
    u_{22} + u_{13} + u_{02} + u_{11} - 4u_{12} &= \num{4,5}^2 \cdot (1 \cdot \num{4,5}  + 2\cdot \num{4,5} ) \\
    100 + u_{13} + 100 + u_{11} - 4u_{12} &= \num{273,375}  \\
    u_{11} - 4u_{12} + u_{13} &= \num{73,375} 
  .\end{align*}
  To find an expression for $u_{13}$ we use the directional derivative as:
  \begin{align*}
    \num{120,25} = u_n(P_{12}) = \frac{\partial u_{12}}{\partial y} &\approx \frac{u_{13} - u_{11}}{9} \\
    u_{13} &= u_{11} + \num{1082,25} 
  .\end{align*}
  By insertion we get:
  \begin{align*}
    u_{11} - 4u_{12} + u_{11} + \num{1082,25}  &= \num{73,375}  \\
    2u_{11} - 4u_{12} &= \num{-1008,875} 
  .\end{align*}
\end{derivation}

\finalresult{
\begin{aligned}
  -4u_{11} + u_{12} &= \num{-117,75}  \\
  2u_{11} - 4u_{12} &= \num{-1008,875} 
.\end{aligned}
}

\begin{problemwithimage}{p13_3}{13.3}
  Consider the two-dimensional Laplace equation
  \begin{equation*}
    u_{x x}(x,y) + u_{yy}(x,y) = 0
  \end{equation*}
  on the region $R$ with boundary conditions as depicted on the figure.
  Find the approximate solution for mesh size $h = \num{1,5}$.
\end{problemwithimage}

\begin{derivation}
  We know the boundary values:
  \begin{align*}
    u_{10} &= \num{1,5}, & u_{01} &= \num{1,5}, & u(A) &= 1, & u(B) &= 1
  .\end{align*}
  We now approximate $\nabla^2 u (P_{11})$ using $a = b = \num{0,5} / \num{1,5}  = 1 / 3$ in the formula from the Lecture (p. 5):
  \begin{align*}
    \nabla^2u(P_{11}) &\approx \frac{2}{\num{1,5}^2} \left( \frac{9}{4}u(A) + \frac{9}{4}u(B) + \frac{3}{4}u(P_{01}) + \frac{3}{4} u(P_{10}) - 6u (P_{11}) \right) \\
    &= \frac{2}{\num{1,5}^2} \left( \frac{9}{4} + \frac{9}{4} + \frac{3}{4} \cdot \num{1,5}  + \frac{3}{4}\cdot \num{1,5} - 6u_{11} \right)
  .\end{align*}

  Hence, our difference equation is
  \begin{align*}
    \nabla^2 u(P_{11}) &= 0 \\
    \frac{2}{\num{1,5}^2} \left( \frac{9}{4} + \frac{9}{4} + \frac{3}{4} \cdot \num{1,5}  + \frac{3}{4}\cdot \num{1,5} - 6u_{11} \right) &= 0 \\
    \num{6,75}  - 6 u_{11} &= 0 \\
    u_{11} &\approx \num{1,125} 
  .\end{align*}
\end{derivation}

\finalresult{
u_{11} \approx \num{1,125}
}

\begin{problemwithimage}{p13_3}{13.4}
  Consider the two-dimensional Poisson equation
  \begin{equation*}
    u_{x x}(x,y) + u_{yy}(x,y) = x + y
  \end{equation*}
  on the region $R$ with boundary conditions as depicted on the figure.
  Find the approximate solution for mesh size $h = \num{1,5}$.
\end{problemwithimage}

\begin{derivation}
  In this case, we have the same boundary values:
  \begin{align*}
    u_{10} &= \num{1,5} , & u_{01} &= \num{1,5} , & u(A) &= 1, & u(B) &= 1
  \end{align*}
  and the same approximation of $\nabla^2 u(P_{11}) $ as above:
  \begin{equation*}
    \nabla ^2 u(P_{11}) \approx \frac{2}{\num{1,5}^2} \left( \frac{9}{4} + \frac{9}{4} + \frac{3}{4} \cdot \num{1,5}  + \frac{3}{4}\cdot \num{1,5} - 6u_{11} \right) 
  .\end{equation*}
  
  In this case, however, our difference equation becomes:
  \begin{align*}
    \nabla^2 u(P_{11}) &= h + h \\
    \num{6,75}  - 6u_{11} &= 3 \cdot \frac{\num{1,5}^2}{2}\\
    u_{11} &\approx \num{0,5625} 
  .\end{align*}
\end{derivation}

\finalresult{
  u_{11} \approx \num{0,5625}
}
